\documentclass[sigconf]{acmart}

\AtBeginDocument{%
  \providecommand\BibTeX{{%
    Bib\TeX}}}

\acmConference[ITSEC]{IT-Security Seminar}{WiSe 2023/24}{Heidelberg}
\begin{document}
\title{PQC Wireguard as a new VPN}
\subtitle{Wireguard + Rosenpass}

\author{Philip Kastura-Sahl}
\email{philip.kastura-sahl@stud.uni-heidelberg.de}
\affiliation{%
  \institution{Ruprecht-Karls-Universität Heidelberg}
  \city{Heidelberg}
  \country{Germany}
}

\renewcommand{\shortauthors}{Philip Kastura-Sahl}

\begin{abstract}
  With the advance of quantum computers security researchers face the challenge of securing their data against \textit{"Harvest now, decrypt later"-attacks} that could become available in the future. Shor's algorithm\cite{Shor_1997}, when implemented on a sufficiently powerfull quantum computer, could be used to break many public-key cryptography schemes. Rosenpass aims to provide a \textit{post-quantum-secure authenticated key exchange} that works in a \textit{hybrid post-quantum security scheme} together with Wireguard. Due to the ongoing development of Post-quantum-Ciphers Rosenpass is implemented to provide hybrid post-quantum security in tandem with Wireguard.\cite{rosenpass-about}
\end{abstract}

%% TODO: refresh eventual tag changes
\begin{CCSXML}
<ccs2012>
   <concept>
       <concept_id>10002978.10002979.10002980</concept_id>
       <concept_desc>Security and privacy~Key management</concept_desc>
       <concept_significance>300</concept_significance>
       </concept>
   <concept>
       <concept_id>10002978.10002979.10002981.10011602</concept_id>
       <concept_desc>Security and privacy~Digital signatures</concept_desc>
       <concept_significance>500</concept_significance>
       </concept>
   <concept>
       <concept_id>10002978.10002979.10002981.10011745</concept_id>
       <concept_desc>Security and privacy~Public key encryption</concept_desc>
       <concept_significance>500</concept_significance>
       </concept>
 </ccs2012>
\end{CCSXML}

\ccsdesc[300]{Security and privacy~Key management}
\ccsdesc[500]{Security and privacy~Digital signatures}
\ccsdesc[500]{Security and privacy~Public key encryption}

\keywords{Post-quantum cryptography, Key-Encapsulation Method, Rosenpass, Wireguard, Crystals-Kyber, NIST}

\begin{teaserfigure}
  \centering
  \includegraphics[width=\textwidth]{graphics/rosenpass+wireguard.pdf}
\end{teaserfigure}

\received[Talk]{18 January 2024}
\received[Paper]{31 March 2024}

\maketitle

%%%%%%%%%%%%%%%%%%%%%%%%%%%%%%%%%%%%%%%%%%%%%%%%%%%%%%%%%%%%%%%%%%%%%%%%%%%%%%%%%%%%%%%%%%%%%%%%%%%%%%%%%%%%%%%%%%%%%%%%%%%%%%%%%%%%%%%%%%%%%%%%%%%%%%%%%%%%%%%%%%%%
% Begin Document
%%%%%%%%%%%%%%%%%%%%%%%%%%%%%%%%%%%%%%%%%%%%%%%%%%%%%%%%%%%%%%%%%%%%%%%%%%%%%%%%%%%%%%%%%%%%%%%%%%%%%%%%%%%%%%%%%%%%%%%%%%%%%%%%%%%%%%%%%%%%%%%%%%%%%%%%%%%%%%%%%%%%

\section{Introduction}
Shor's Algorithm, developed in 1994 by Peter Shor, is one of the few known algorithms that could be used to used to break existing public key cryptography such as RSA and some Diffie-Hellman variants. With the rise of ever more powerful quantum computers security researchers are facing the possibility of traditional \textit{pre-quantum cryptography} becoming insecure. The Open Quantumn Safe Project\cite{open-quantum-safe} \& NIST\footnote{NIST -- National Institute of Standards and Technology} aim to improve and standardize \textit{post-quantum cryptography} to the point of it being able to supersede existing schemes that would be deemed insecure using sufficiently powerful quantumn computers. \\
PQ-Wireguard\cite{cryptoeprint:2020/379} first modified Wireguard and implemented new PQ\footnote{pq -- post-quantum} forward secrecy \& pq Authentication by replacing the Diffie-Hellman handshake with a key-encapsulation mechanism\footnote{KEM -- key-encapsulation mechanism}.

\subsection{Wireguard \& Rosenpass}
% TLDR: What is Wireguard; What is Rosenpass
Wireguard is a modern VPN that intends to be more performant that OpenVPN, whilst being able to run on a variety of computers and platforms\cite{Wireguard}. Rosenpass functions as an Addon to Wireguard, adding PQC\footnote{PQC -- post-quantum cryptography} whilst keeping the established security intact. Whilst some Wireguard implementations are considered complete\cite{wireguard-repos}, Rosenpass is still under development\cite{rosenpass}. Apart from implementing PQC, Rosenpass also implements a statless responder to protect against replay attacks of the first protocol message.

\subsection{Problem Definition}
% TLDR: Explain pq Computing
Even though Quantumn Computers don't yet have a sufficient number of qubits, data could be harvested already, making it vulnerable to decryption attacks in the future, a surveillance strategy known as \textit{Harvest now, decrypt later}. Yet PQC has not yet reached the point of adoption to be replacing pre-quantum cryptography, as seen in the Post-Quantum Cryptography Standardization effort where algorithms such as SIDH have been eliminated due to major security flaws in their design\cite{cryptoeprint:2022/975}. Therefore Rosenpass has been designed to implement Hybrid PQC\cite{cryptoeprint:2022/1225}, Rosenpass providing encryption deemed the most qost-quantum resistant and working together with the known to be pre-quantumn secure Wireguard. The handshake between Wireguard and Rosenpass as well as the individual connections between Rosenpass and Wireguard Clients and Server needed to be designed in such a way that neither one of the ciphers could be broken alone so that a potential attacker would need to always attack both ciphers.


\section{Ciphers}
% TLDR: NIST Difficulties with Ciphers
The Rosenpass developers have chosen a combination of Classic McEliece and CRYSTALS-Kyber, notably finalists of the NIST Post-Quantum Cryptography Standardization effort\cite{pqc-standardization}. While symmetric ciphers are generally considered to be secure, not all asymmetric ciphers are equally vulnerable to Shor's algorithm.

\subsection{KEM vs. NIKE}
% TLDR: What is KEM; What is NIKE; Why do we use NIKE?
Established asymmetric ciphers \textit{Non-Interactive Key Exchange}, as the name implies, this enables two parties to know each others public keys and agree on a symmetric shared key without requiring any interaction\cite{cryptoeprint:2012/732}. While a NIKE\footnote{NIKE -- Non-Interactive Key Exchange} is not inherently insecure against attacks by quantum computers, PQC usually relies on a \textit{Key encapsulation mechanism}. KEMs\footnote{KEM -- Key encapsulation mechanism} are considered a more viable solution, because they are able to provide both pq-resistance as well as today's security guarantees. Hybrid PQC can be implemented using KEMs by using a combiner, which enables the use of multiple algorithms while keeping the scheme secure as long as one of the algorithms remains secure\cite{cryptoeprint:2012/732}.

%\subsection{Classic McElice}
%% TLDR: Explain Classic McElice
% Classic McEliece uses \textit{goppa codes}

\subsection{Kyper}
% TLDR: Explain Kyper
\textit{Lattice-based ciphers} such as CRYSTALS-Kyber


\section{Implementation}
% TLDR: Rough explainer how Rosenpass is developed along Wireguard

\subsection{Proof}
% TLDR: State of Mathematical Proof; Proof only written using Pseudocode?

\subsection{Key Exchange}
% TLDR: Explain Key Exchange

\subsection{Wireguard-Rosenpass Handshake}
% TLDR: Explain Handshake


\section{Conclusion}
% TLDR: Conclude with Usability and Provability; Security of KEMs???

%%%%%%%%%%%%%%%%%%%%%%%%%%%%%%%%%%%%%%%%%%%%%%%%%%%%%%%%%%%%%%%%%%%%%%%%%%%%%%%%%%%%%%%%%%%%%%%%%%%%%%%%%%%%%%%%%%%%%%%%%%%%%%%%%%%%%%%%%%%%%%%%%%%%%%%%%%%%%%%%%%%%

\begin{acks}
I would like to express my sincere gratitude to the following individuals and teams who have contributed to the completion of this document:

\begin{itemize}
  \item Prof. Dr. Heuveline \& Mr. Machmeier for their guidance and support throughout the research and writing process.
  \item The Wireguard \& Rosenpass teams for their respective Projects, their novice-friendly documentation, and beginner-friendly explainers.
\end{itemize}

I am truly thankful for their dedication and collaboration, which has significantly enriched the quality of this work.
\end{acks}

%% TODO: refresh accessed citations
\bibliographystyle{ACM-Reference-Format}
\bibliography{sources}

%% TODO: Use Appendix if there are reused Definitions
%\appendix

%%%%%%%%%%%%%%%%%%%%%%%%%%%%%%%%%%%%%%%%%%%%%%%%%%%%%%%%%%%%%%%%%%%%%%%%%%%%%%%%%%%%%%%%%%%%%%%%%%%%%%%%%%%%%%%%%%%%%%%%%%%%%%%%%%%%%%%%%%%%%%%%%%%%%%%%%%%%%%%%%%%%
% End Document
%%%%%%%%%%%%%%%%%%%%%%%%%%%%%%%%%%%%%%%%%%%%%%%%%%%%%%%%%%%%%%%%%%%%%%%%%%%%%%%%%%%%%%%%%%%%%%%%%%%%%%%%%%%%%%%%%%%%%%%%%%%%%%%%%%%%%%%%%%%%%%%%%%%%%%%%%%%%%%%%%%%%

\end{document}
\endinput
